
\chapter{\textcolor{red}{Electric Motors}}
Electric motors are fundamental components in countless systems, converting electrical energy into mechanical motion. In this section, we will explore the main types of AC and DC motors, examining their operating principles and key characteristics. We will also look at the signals and driver circuits that enable these motors to function effectively.

\Mynote{Add more general intro}
\section{Electromagnetic Principles for Motor Applications}
Most motors rely on the principle that a current-carrying conductor placed in a magnetic field experiences a force \cite[see][2]{Hughes:2006}. As shown in Figure \ref {fig:lorentz}, an interaction occurs between the current-generated magnetic field and the permanent magnetic field, which is expressed through the so-called Lorentz force, named after the Dutch physicist Hendrik Antoon Lorentz \cite[see][60]{Sturm:2018}. This force can be observed directly by placing a bar magnet near a wire through which current flows. When the current flows through the conductor, it interacts with the magnetic field, creating a force that acts on the conductor itself.  Experiments have shown that the magnitude of this force depends directly on the amount of current in the wire as well as the strength of the magnetic field. Furthermore, the force is at its maximum when the magnetic field is oriented perpendicular to the conductor \cite[see][2]{Hughes:2006}.

\begin{figure}[H]
	\centering
	\includegraphics[width=0.7\linewidth]{../../../Images/Syde/MotorControl/Lorentz}
	\caption[ Mechanical force produced on a current-carrying wire in a magnetic Weld
	]{ Mechanical force produced on a current-carrying wire in a magnetic Weld
		 \cite[see][2]{Hughes:2006}}
	\label{fig:lorentz}
\end{figure}

The force $F[N]$ acting on a wire of length $l[m]$, carrying a current $I[A]$ and exposed to a uniform magnetic flux density $B[T]$, is given by a straightforward expression \ref{eq:lorentzForce}:

\begin{equation}
	F=BIl
	\label{eq:lorentzForce}
\end{equation} 

In most motors, the working magnetic field is generated by coils of wire carrying current. To understand this, the simplest scenario is considered: the magnetic field surrounding an isolated long straight wire that carries a steady current as depicted in Figure \ref{fig:FluxOnOneWire}. In the figure, the plus sign indicates current flowing into the page, while a dot represents current coming out. This convention can be visualized by imagining an arrow or dart, with the cross showing the rear of the shaft and the dot as the tip. Around this wire, the magnetic flux lines form concentric circles, with the field strength greatest near the wire \cite[see][7-8]{Hughes:2006}.

\begin{figure}[H]
	\centering
	\includegraphics[width=0.7\linewidth]{../../../Images/Syde/MotorControl/FluxOnOneWire}
	\caption[ MMagnetic flux lines produced by a straight, current-carrying wire]{Magnetic flux lines produced by a straight, current-carrying wire \cite[see][8]{Hughes:2006}}
	\label{fig:FluxOnOneWire}
\end{figure}

Now, if a parallel “go” and “return” circuit is considered, the magnetic field is the result of superimposing the fields generated by the positive current in one conductor and the negative current in the return conductor, as shown in Figure \ref{fig:FluxOnTwoWires}. This superimposition leads to an increased magnetic field in the region between the conductors and a reduced field outside this region. While this configuration applies exactly to an infinitely long pair of conductors, it also describes the field generated by a single turn of wire—regardless of whether the turn is rectangular, square, or round in shape. This concept helps build a clearer understanding of the magnetic fields created by the coils typically used in motors, which often consist of many turns \cite[see][8-9]{Hughes:2006}.

\begin{figure}[H]
	\centering
	\includegraphics[width=0.7\linewidth]{../../../Images/Syde/MotorControl/FluxOnTwoWires}
	\caption[ Magnetic flux lines produced by current in a parallel go and return circuit]{  Magnetic flux lines produced by current in a parallel go and return circuit \cite[see][8]{Hughes:2006}}
	\label{fig:FluxOnTwoWires}
\end{figure}

In Figure \ref{fig:coil}, the coil is displayed on the left, while the flux pattern it produces is shown on the right. Each turn of the coil contributes its own magnetic field pattern. When these individual fields are combined, the field inside the coil becomes significantly stronger. Between each pair of flux lines, there is an equal amount of flux. As the flux lines spread outward from the coil, the flux density outside the coil is much lower than inside. For instance, if the distance b is four times the distance a, the flux density $B_b$ will be a quarter of $B_a$ \cite[see][9]{Hughes:2006}.

\begin{figure}[H]
	\centering
	\includegraphics[width=0.7\linewidth]{../../../Images/Syde/MotorControl/coil}
	\caption[Multi-turn cylindrical coil and pattern of magnetic flux produced by current in the coil.]{  Multi-turn cylindrical coil and pattern of magnetic flux produced by current in the coil. \cite[see][9]{Hughes:2006}}
	\label{fig:coil}
\end{figure}

\paragraph{MMF and EMF} To describe a coil’s ability to generate magnetic flux, we use the term Magnetomotive force (MMF). MMF is simply the product of the number of turns in the coil and the current flowing through it, measured in ampere-turns. Consequently, one straightforward way to increase the magnetic flux density or MMF in a coil is to either increase the current or add more turns of wire, since both factors are directly proportional to the total flux and thus to the flux density throughout the system. This relationship is similar to what we observe in electrical circuits, where the current is directly proportional to the applied voltage. In electrical systems, this relationship is described by Ohm’s law, which states that the current depends on both the Electromotive force (EMF) and the circuit resistance: to increase the current, the resistance must be reduced \ref{eq:EMF}. 

\begin{equation}
	Current = \frac{EMF}{Resistance}		i.e.		I = \frac{V}{R}
	\label{eq:EMF}
\end{equation} 

In magnetic circuits, we encounter a parallel concept known as reluctance ($\mathcal{R}$). Reluctance measures how difficult it is for magnetic flux to complete its path, much like resistance quantifies how much opposition electric current faces in an electrical circuit. The magnetic equivalent of Ohm’s law states that for a given MMF, the magnetic flux depends on the circuit’s reluctance.

\begin{equation}
	Flux = \frac{MMF}{Reluctance}		i.e.		\phi = \frac{NI}{\mathcal{R}}
	\label{eq:MMF}
\end{equation} 

\paragraph{Iron Core}From this, it follows that to achieve a higher flux for a given MMF, the reluctance of the magnetic path must be reduced. In the case of a coil, this involves replacing as much of the air gap as possible with a material that better conducts magnetic flux, since air has poor magnetic properties and high reluctance. This substitution lowers the reluctance and allows for a higher magnetic flux with the same MMF. The preferred material for this purpose is high-quality magnetic steel, commonly referred to as “iron.” This material offers significant advantages, which are illustrated in Figure \ref{fig:CoilCore} \cite[see][9-10]{Hughes:2006}.

\begin{figure}[H]
	\centering
	\includegraphics[width=0.7\linewidth]{../../../Images/Syde/MotorControl/CoilCore}
	\caption[Flux lines inside low-reluctance magnetic circuit with air-gap]{ Flux lines inside low-reluctance magnetic circuit with air-gap \cite[see][11]{Hughes:2006}}
	\label{fig:CoilCore}
\end{figure}


Firstly, the reluctance of the iron paths is significantly lower than that of the air paths they replace, so the total magnetic flux produced for a given MMF is considerably greater. Secondly, nearly all of the flux is confined within the iron, rather than spreading out into the surrounding air. This confinement allows to shape the iron parts of the magnetic circuit, as illustrated in Figure \ref{fig:CoilCore}, to guide the flux precisely where it is needed. Finally, within the iron, the flux density remains uniform across the entire cross-section. The extremely low reluctance of the iron ensures there is no noticeable tendency for the flux to crowd to one side or another. 

\paragraph{The Air-gap} Because it is physically impossible to place current-carrying conductors inside the iron itself, an air gap is introduced into the magnetic circuit. If the air gap is relatively small, as in the case of motors, the magnetic flux will jump across the gap as depicted in Figure \ref{fig:CoilCore}, with minimal spreading into the surrounding air. This means that most of the flux lines traverse the air gap directly, and the flux density in this region remains as high as it is within the iron \cite[see][11-12]{Hughes:2006}. 

\begin{figure}[H]
	\centering
	\includegraphics[width=0.7\linewidth]{../../../Images/Syde/MotorControl/saturation}
	\caption[Sketch showing how the effective reluctance of iron increases rapidly as the
	flux density approaches saturation]{ Sketch showing how the effective reluctance of iron increases rapidly as the flux density approaches saturation \cite[see][15]{Hughes:2006}}
	\label{fig:saturation}
\end{figure}

\paragraph{Saturation} As long as the magnetic flux density remains below approximately 1.6 to 1.8 Tesla, the reluctance of the iron parts of the circuit remains minimal. However, when the flux density increases beyond this range, a significant portion of the magnetomotive force (MMF) is consumed in driving the flux through the iron, leaving less MMF to drive the flux across the air gap. This phenomenon, where the reluctance of the iron begins to noticeably increase, is referred to as the onset of magnetic saturation \ref{fig:saturation}. Once the iron starts to saturate, further increases in MMF or reductions in the cross-sectional area of the iron lead only to marginal increases in flux density, which typically levels off around 2 tesla. To avoid the detrimental effects of saturation, the dimensions of the iron components in the magnetic circuit are generally designed so that the flux density does not exceed about 1.5 tesla. At this level, the reluctance of the iron remains low compared to the reluctance of the air gap, ensuring efficient flux transfer \cite[see][14-15]{Hughes:2006}.
 
\paragraph{Simple DC motor}The simplest form of a DC motor features a two-pole design, with an armature that contains just two conductors forming a single coil \ref{fig:simpleDC}. A commutator with two segments and carbon brushes supplies the armature with current from an external direct voltage source. Simultaneously, another direct voltage source supplies current to the field winding, establishing the magnetic flux between the two poles. The armature current flows through the winding and creates additional magnetic fields around the conductors. 

\begin{figure}[H]
	\centering
	\includegraphics[width=0.7\linewidth]{../../../Images/Syde/MotorControl/simpleDC}
	\caption[Simplest form of dc motor showing armature and stator windings]{Simplest form of dc motor showing armature and stator windings \cite[see][269]{Clifford:2010}}
	\label{fig:simpleDC}
\end{figure}

These subsidiary fields interact with the main field, resulting in a combined field depicted in Figure \ref{fig:ForceDC}. This interaction produces forces on the conductors that generate a torque, causing the armature to rotate \cite[see][269]{Clifford:2010}.

\begin{figure}[H]
	\centering
	\includegraphics[width=0.4\linewidth]{../../../Images/Syde/MotorControl/ForceDC}
	\caption[Generation of forces on a pair of conductors carrying current]{Generation of forces on a pair of conductors carrying current \cite[see][270]{Clifford:2010}}
	\label{fig:ForceDC}
\end{figure}

To maintain rotation in the same direction, the commutator and carbon brushes reverse the current in the armature conductors every time the coil crosses the plane perpendicular to the main magnetic field. Figure \ref{fig:ArmatureOrientations} shows how the current in the armature conductors alternates direction as the armature rotates, ensuring continuous torque in a consistent rotational direction.

\begin{figure}[H]
	\centering
	\includegraphics[width=0.9\linewidth]{../../../Images/Syde/MotorControl/ArmatureOrientations}
	\caption[Forces on armature windings at different armature orientations]{Forces on armature windings at different armature orientations \cite[see][270]{Clifford:2010}}
	\label{fig:ArmatureOrientations}
\end{figure}

For analysis purposes, the armature winding can be represented by an equivalent circuit \ref{fig:CircuitDC} that includes the induced EMF and the armature resistance. Typically, the armature resistance is quite low. The armature terminal voltage in this model is given by \ref{eq:TermianalVoltage} \cite[see][14-15]{Clifford:2010}:

\begin{equation}
	V_a = E + I_aR_a
	\label{eq:TermianalVoltage}
\end{equation}

\begin{figure}[H]
	\centering
	\includegraphics[width=0.5\linewidth]{../../../Images/Syde/MotorControl/CircuitDC}
	\caption[Simplest form of dc motor showing armature and stator windings]{Simplest form of dc motor showing armature and stator windings \cite[see][269]{Clifford:2010}}
	\label{fig:CircuitDC}
\end{figure}


\paragraph{Motor Evolution} The progression from a basic C-core configuration to the motor geometry used in a typical DC motor is illustrated in \ref{fig:motorEvolution}. In the first stage of this evolution (\ref{fig:motorEvolution}, left), the original single air gap of length $g$ is divided into two gaps of length $\frac{g}{2}$ each, reflecting the requirement for the rotor to rotate. At the same time, the single magnetizing coil is split into two coils to maintain symmetry. It is worth noting that the position of the magnetizing coil can be adjusted within the magnetic circuit, much like a battery can be placed at various points within an electrical circuit.
In the next stage (\ref{fig:motorEvolution}, center), the single magnetic path is divided into two parallel paths, each with half the original cross-sectional area and carrying half of the total flux. Finally, in the configuration shown in \ref{fig:motorEvolution} on the right, the flux paths and the pole faces are curved to fit the rotor geometry. The coil is wound in several layers to utilize the available space efficiently, but as previously discussed, this layering does not affect the MMF. The air gap remains small, ensuring that the flux crosses the gap radially to reach the rotor \cite[see][16]{Hughes:2006}.

\begin{figure}[H]
	\centering
	\includegraphics[width=0.9\linewidth]{../../../Images/Syde/MotorControl/motorEvolution}
	\caption[Evolution of d.c. motor magnetic circuit from gapped C-core]{Evolution of d.c. motor magnetic circuit from gapped C-core \cite[see][16]{Hughes:2006}}
	\label{fig:motorEvolution}
\end{figure}

\paragraph{Torque} To create rotational motion, a set of conductors is arranged and fixed to the rotor, as illustrated in Figure \ref{fig:motorConductors}. In this setup, the conductors located under the north pole (at the top of Figure \ref{fig:motorConductors}) carry positive current, flowing into the page, while those under the south pole carry negative current. The tangential electromagnetic force (often referred to as the "BIl" force, see Equation \label{eq:lorentzForce}) acts to the left on the positive conductors and to the right on the negative ones.

\begin{figure}[H]
	\centering
	\includegraphics[width=0.2\linewidth]{../../../Images/Syde/MotorControl/motorConductors}
	\caption[Current-carrying conductors on rotor, positioned to maximise torque.]{Current-carrying conductors on rotor, positioned to maximise torque. \cite[see][17]{Hughes:2006}}
	\label{fig:motorConductors}
\end{figure}

This generates a net torque on the rotor, causing it to rotate. From this, it follows that the total tangential force $F$ acting on the rotor depends on the flux density produced by the field winding, the number of conductors on the rotor, the current flowing in each conductor, and the length of the rotor. The torque $T$ produced as a result of this tangential force also depends on the radius $r$of the rotor \ref{eq:torque} \cite[see][16-19]{Hughes:2006}.

\begin{equation}
	T = Fr
	\label{eq:torque}
\end{equation} 

\paragraph{Slotting} Initially, rotor conductors were mounted on the surface of the rotor, requiring a large air gap for clearance and secure attachment. This led to high reluctance in the magnetic circuit, demanding bulky field windings with high current. To reduce the air gap and improve efficiency, conductors were soon recessed into slots machined along the rotor \ref{fig:slotted}. These “slotted” machines not only reduced the reluctance and air gap size, but also provided better mechanical support for the conductors.

\begin{figure}[H]
	\centering
	\includegraphics[width=0.3\linewidth]{../../../Images/Syde/MotorControl/slotted}
	\caption[Influence on flux paths when the rotor is slotted to accommodate conductors]{Influence on flux paths when the rotor is slotted to accommodate conductors \cite[see][19]{Hughes:2006}}
	\label{fig:slotted}
\end{figure}

As conductors were sunk deeper, it became apparent that although the flux largely bypassed the conductors in the slots, the torque produced was the same as if the conductors were in the main flux path. This is because most of the tangential force now acts on the rotor teeth, which are sturdy and well-suited to transmit torque to the shaft. An added benefit is that the conductor insulation no longer needs to withstand tangential forces, simplifying design requirements.

To maximize torque, high current densities (typically $\frac{2-8A}{mm^2}$) and larger cross-sectional slot areas are used to fit as much copper as possible. However, making the slots too wide narrows the teeth, risking magnetic saturation. Similarly, deepening the slots too much can deplete the central rotor iron, which also leads to saturation issues \cite[see][19-20]{Hughes:2006}.

\begin{figure}[H]
	\centering
	\includegraphics[width=0.5\linewidth]{../../../Images/Syde/MotorControl/BrushedDC}
	\caption[Conventional (brushed) D.C. motor]{Conventional (brushed) D.C. motor \cite[see][83]{Hughes:2006}}
	\label{fig:BrushedDC}
\end{figure}


\paragraph{Brushed DC Motor} The culmination of these Electromagnetic Principles can be seen in the classic design of a conventional (brushed) DC motor. In these motors, the rotor teeth and slots ensure efficient torque transmission, while the brush and commutator system manages the continuous reversal of current to sustain rotation. With this foundational design established, it is instructive to examine how different types of motors incorporate these principles in unique ways to meet diverse application requirements.



\section{Motor Types and Their Characteristics}
Motors are typically categorized by their type and the kind of electrical power they use. In alternating current (AC) machines, there are two main types. The first is the polyphase synchronous motor, where the rotor’s magnetic field is generated by either direct current supplied through slip rings or by permanent magnets mounted directly on the rotor. The second type is the induction motor, which can be either single-phase or polyphase. In this type, the rotor’s magnetic field is created by induction effects, rather than direct electrical connection. Direct current (DC) motors, on the other hand, are classified based on how the magnetic field is connected in the circuit, such as series, shunt, or compound configurations.

Despite these differences, all motors share two fundamental requirements. First, to generate a continuous torque, the magnetic fields in the stator and rotor must remain stationary with respect to one another. In AC motors, this is achieved by both fields rotating in sync. In DC motors, both fields remain stationary relative to the motor’s frame. Second, for a motor with a given rotor length, air gap diameter, and rotational speed, there is a maximum achievable power output. This limitation is dictated by the maximum magnetic and electrical loading allowed by the materials used. The maximum magnetic loading is determined by how much the steel in the magnetic circuit can be magnetized without saturating. The flux density in the air gap is limited to ensure the armature teeth do not become oversaturated. The electrical loading, meanwhile, depends on the current density in the armature windings. This current density must be managed carefully to avoid excessive heating, which is linked to how well the system can dissipate heat and to the temperature limits of the insulation used around the conductors \cite[see][26]{Toliyat:2004}.

\subsection{Direct Current Motors}
Brushed DC motors stand out for their simple design and low initial cost, offering reliable performance along with straightforward control of speed and torque. Brushless DC motors have since evolved as power electronic technologies advanced, replacing the mechanical commutation of brushes with electronic switching. This update brings several advantages: brushless motors are more efficient, last longer, and require less maintenance. They also reduce noise and eliminate the sparks and electromagnetic interference that often accompany brushed commutators. Today, DC motors continue to play a crucial role in applications requiring precise speed control over a wide range around the rated speed—such as steel mills, mining operations, and electric trains. However, these motors also face challenges. Brushes and commutators, essential for delivering power to the rotor, introduce extra complexity and wear. Commutators can generate power ripples and limit rotor speed, while brushes add frictional losses and can generate unwanted radiofrequency interference. Over time, brush wear also creates carbon dust inside the motor, which can lead to reliability and performance issues. As a result, maintaining the brush–commutator interface is critical, but it also adds to the operational complexity of these motors. Additionally, the brushes themselves contribute to increased acoustic noise \cite[see][39]{Tong:2014}. 

There are several different types of DC motors.

\subsubsection{Shunt Motors}
In shunt motors \ref{CircuitShunt}, the rotor and stator windings are connected in parallel, which keeps the currents in these circuits independent. This design ensures excellent speed regulation, making these motors reliable for steady operation at varying loads \cite[see][38]{Tong:2014}.

\begin{figure}[H]
	\centering
	\includegraphics[width=0.5\linewidth]{../../../Images/Syde/MotorControl/CircuitShunt}
	\caption[Circuit for a Shunt motor ]{Circuit for a Shunt motor \cite[see][273]{Clifford:2010}}
	\label{fig:CircuitShunt}
\end{figure}

\subsubsection{Separately Excited Motors}
These motors use separate power sources for the rotor and stator windings \ref{fig:CircuitSepExcited}. Common in industrial settings, their independent control of the magnetic fields provides precise speed and position control, making them ideal for automation and feedback systems \cite[see][38]{Tong:2014}.

\begin{figure}[H]
	\centering
	\includegraphics[width=0.7\linewidth]{../../../Images/Syde/MotorControl/CircuitSepExcited}
	\caption[Circuit for a Separately Excited motor ]{Circuit for a Separately Excited motor \cite[see][273]{Clifford:2010}}
	\label{fig:CircuitSepExcited}
\end{figure}

\subsubsection{Series Motors}
Series motors \ref{fig:CircuitSeries} have the rotor and stator windings connected in series. This arrangement produces torque that scales with the square of the current, resulting in very high starting torque. Consequently, they are used in applications like car starter motors and older lift systems \cite[see][38]{Tong:2014}.

\begin{figure}[H]
	\centering
	\includegraphics[width=0.3\linewidth]{../../../Images/Syde/MotorControl/CircuitSeries}
	\caption[Circuit for a Series motor ]{Circuit for a Series motor \cite[see][273]{Clifford:2010}}
	\label{fig:CircuitSeries}
\end{figure}

\subsubsection{Compound Motors}
Compound motors come in two main configurations: the long-shunt motor and the short-shunt motor. Both designs feature two separate field windings—the shunt-field winding and the series-field winding. In the long-shunt motor, the shunt-field winding is connected across both the armature and the series-field winding, whereas in the short-shunt motor, it is connected only across the armature\ref{fig:CircuitCompound}. In both configurations, the shunt-field winding generates a stronger magnetic field than the series-field winding. Typically, the windings are connected so that the magnetic fields they produce combine to reinforce one another. This configuration is known as cumulative compounding \cite[see][282]{Clifford:2010}. Combining the series and shunt designs, compound motors offer a balance of strong torque and stable speed.\cite[see][38]{Tong:2014}.

\begin{figure}[H]
	\centering
	\includegraphics[width=0.9\linewidth]{../../../Images/Syde/MotorControl/CircuitCompound}
	\caption[Long-shunt compound motor and Short-shunt compound motor]{Long-shunt compound motor and Short-shunt compound motor \cite[see][273]{Clifford:2010}}
	\label{fig:CircuitCompound}
\end{figure}

\subsubsection{DC Permanent Magnet (PM) Motors}
In these motors, permanent magnets replace the field windings \ref{fig:PermanentDC}, simplifying the design and improving efficiency. Mechanical commutation is managed by a series of carbon brushes and commutator segments, ensuring continuous rotation \cite[see][38]{Tong:2014}.

\begin{figure}[H]
	\centering
	\includegraphics[width=0.4\linewidth]{../../../Images/Syde/MotorControl/PermanentDC}
	\caption[2-pole permanent magnet motor]{2-pole permanent magnet \cite[see][85]{Hughes:2006}}
	\label{fig:PermanentDC}
\end{figure}

\subsubsection{Universal Motors}
The series commutator motor, more commonly known as the universal motor, is widely used in portable power tools, food mixers, vacuum cleaners, and similar devices. Despite being originally designed for DC power, these motors are can be operated on AC supplies as well. In a series motor, the field flux is produced by the same current that flows through the armature. As a result, when the current reverses direction as it does in an AC supply, the magnetic flux reverses as well. This ensures that the torque generated by the interaction between the armature current and the magnetic field remains in the same direction, regardless of the current’s polarity. For example, when powered by a 50\,Hz AC supply, the current reverses every 10 milliseconds, producing torque peaks 100 times per second.

\begin{figure}[H]
	\centering
	\includegraphics[width=0.5\linewidth]{../../../Images/Syde/MotorControl/UniversalMotor}
	\caption[Universal motor principles]{Universal motor principles \cite[see][26]{Toliyat:2004}}
	\label{fig:UniversalMotor}
\end{figure}

To handle the demands of AC operation, universal motors are built with fully laminated construction to minimize eddy current losses caused by the alternating flux in the magnetic circuit. These motors are designed for high-speed operation, typically reaching 8,000 to 12,000 revolutions per minute at rated voltage. Although commutation and sparking tend to be more problematic when operating on AC rather than DC, these motors can still achieve output powers up to about 1\,kW. Their compact size and high power density make them particularly well suited for small, powerful devices \cite[see][118-119]{Hughes:2006}.



\subsection{Alternating Current motors}
In contrast to DC motors, AC motors rely on alternating current power sources and avoid mechanical components like commutators and brushes entirely. By eliminating these, AC motors achieve higher operational reliability, greater efficiency, and reduced maintenance. Their simpler structure also translates into lower costs and more compact designs. These features make AC motors a popular choice in industrial motion control, where their performance and cost-effectiveness are particularly appealing. Single-phase AC motors are commonly used in residential and commercial applications, while three-phase AC motors are preferred for high-power demands. However, single-phase motors tend to be less efficient than their three-phase counterparts.
Despite their many advantages, AC motors are not without drawbacks. They can be more challenging to control, particularly when precise speed adjustments are needed. Their control systems are often more complex, and they offer lower torque density compared to some DC designs. In addition, AC motors face limitations when operating at low speeds and can suffer from eddy current and hysteresis losses within the stator and rotor cores. These factors, along with higher costs and less precise positioning control, highlight the trade-offs that come with choosing AC motors for specific applications \cite[see][39]{Tong:2014}.

\subsubsection{Polyphase synchronous motors}
AC synchronous motors share similarities with Brushless DC motors. The main difference is the waveform used to drive them: while AC synchronous motors typically use a sinusoidal waveform—either from a direct AC supply or generated via pulse-width modulation—brushless DC motors often rely on rectangular or trapezoidal waveforms. As a result, brushless DC motors develop a rectangular or trapezoidal back EMF, in contrast to the sinusoidal back EMF seen in AC synchronous designs \cite[see][39]{Tong:2014}.
In a synchronous motor, the stator windings are identical to those in an induction motor, creating a rotating magnetic field when connected to a three-phase power supply. However, unlike an induction motor’s rotor, which typically uses a squirrel-cage winding, a synchronous motor has a rotor that features either a DC-excited winding—powered through slip rings—or permanent magnets. These rotor designs are intended to “lock onto” or synchronize with the rotating magnetic field produced by the stator. Once synchronized, the rotor maintains the same speed as the rotating field, ensuring that under constant-frequency operation, the motor’s speed remains constant even as the load varies, provided that the supply frequency itself stays stable \cite[see][39]{Hughes:2006}.

\subsubsection{Induction motors}
\begin{figure}[H]
	\centering
	\includegraphics[width=0.5\linewidth]{../../../Images/Syde/MotorControl/InductionWave}
	\caption[Three-phase ac current waveforms $120^\circ$ apart]{Three-phase ac current waveforms $120^\circ$ apart\cite[see][293]{Clifford:2010}}
	\label{fig:InductionWave}
\end{figure}

\begin{figure}[H]
	\centering
	\includegraphics[width=0.9\linewidth]{../../../Images/Syde/MotorControl/InductionSteps}
	\caption[Sequence of energising the windings of a two-pole, three-phase induction motor]{ Sequence of energising the windings of a two-pole, three-phase induction motor \cite[see][293]{Clifford:2010}}
	\label{fig:InductionSteps}
\end{figure}

Figure \ref{fig:InductionSteps} illustrates the principle of the rotating magnetic field in a three-phase AC induction motor. In this type of motor, three separate stator windings are energized by a balanced three-phase supply \ref{fig:InductionWave}. Each winding is spaced 120 degrees apart within the stator slots. When the three-phase supply is connected, the resulting currents create a magnetic field that rotates at a constant speed in the air gap between the stator and rotor. This rotating magnetic field induces currents in the rotor conductors, which are short-circuited by end rings. These induced currents in the rotor interact with the rotating magnetic field, producing a torque that causes the rotor to turn. The diagrams show how the magnetic field completes a full revolution of the air gap in one cycle of the supply current, continuously driving the rotor forward. This steady rotation and the relative difference in speeds between the rotating field and the rotor (the slip) are fundamental to the operation of induction motors \cite[see][292-293]{Clifford:2010}.

\section{\textcolor{red}{Signal to Motor}}
\subsection{\textcolor{red}{Pulse Width Modulation (PWM)}}
Pulse Width Modulation (PWM) is a technique used to control the power delivered to a load by rapidly alternating the electrical signal between fully on and fully off states. In this approach, the PWM signal resembles a square wave, with the high portion representing the time the power is on (also known as the pulse width PW), and the low portion representing the time it’s off. By adjusting the pulse width—lengthening it to deliver more power or shortening it to deliver less—PWM can effectively manage the amount of power supplied to a device. This switching occurs within a defined cycle called the PWM period T, which includes both the on and off times. The precision of PWM control depends on its resolution, which defines how many discrete steps can fit into one period; higher resolution means finer control. The duty cycle DC, a key term in PWM, quantifies the ratio of the on-time to the total period, expressed as a percentage. For instance, a 0\% duty cycle means the load is always off, while a 100\% duty cycle means it is always on. A duty cycle DC (\%) may be calculated as shown in following Equation \ref{eq:dutyCycle}:

\begin{equation}
	DC = \frac{PW}{T} \times 100\% 
\label{eq:dutyCycle}
\end{equation} 

PWM’s effectiveness lies in how the power to the load is determined by the ratio of the pulse width (the ON time) to the total period of the signal. For example, if the duty cycle is 80\%, it means the load receives power 80\% of the time and is off for the remaining 20\%. A higher duty cycle leads to more power being delivered, while a lower duty cycle reduces it. In this context, PWM frequency is also essential being the inverse of the period \ref{eq:freqPWM}, so shorter periods lead to higher frequencies. This frequency should be high enough that the load (such as a motor or light) cannot respond to the individual on-off pulses but instead perceives an average power level that’s smoother and closer to a continuous signal.

\begin{equation}
	f = \frac{1}{T}
	\label{eq:freqPWM}
\end{equation} 

The PWM waveform typically appears as a series of rectangular pulses \ref{fig:WavePWM}. The ratio of on-time to the total period determines how much energy is transferred to the load. This high-frequency switching is a key advantage because, despite rapidly alternating between on and off states, the energy loss in the switches themselves is kept very low. When the switch is off, there’s essentially no current flow and thus minimal power loss. When it’s on and delivering power, there’s almost no voltage drop across the switch itself, so the power loss is again negligible. This efficient power transfer ensures that PWM remains an excellent choice for many control and power applications.

\begin{figure}[H]
	\centering
	\includegraphics[width=0.7\linewidth]{../../../Images/Syde/MotorControl/WavePWM}
	\caption[A typical waveform of PWM output signal]{A typical waveform of PWM output signal \cite[see][50]{Uzam:2021}}
	\label{fig:WavePWM}
\end{figure}

 In Arduino, PWM signals can be generated using the analogWrite() command, where values of 0, 64, 127, 191, and 255 correspond to the duty cycles of 0\%, 25\%, 50\%, 75\%, and 100\%, respectively. By selecting these values, the motor’s speed can be adjusted to 0\%, 25\%, 50\%, 75\%, or 100\% of its maximum speed.
\begin{figure}[H]
	\centering
	\includegraphics[width=0.7\linewidth]{../../../Images/Syde/MotorControl/PWM}
	\caption[PMW signals generated using analog Write()]{PMW signals generated using analog Write() \cite[see][247]{Biswanath:2023}}
	\label{fig:PWM}
\end{figure}

\subsection{\textcolor{red}{H-Bridge}}
The direction of rotation of a DC motor can be managed using an H-bridge circuit. This circuit consists of four power transistors arranged in an H-like configuration, with the motor at the center (as illustrated in Fig. 9.9). To control the motor’s direction, pairs of transistors in the H-bridge are selectively activated. For instance, turning on the upper left (Q1) and lower right (Q4) transistors while keeping the upper right (Q3) and lower left (Q2) off causes the motor to rotate in one direction (e.g., clockwise). Conversely, turning on the upper right (Q3) and lower left (Q2) transistors while keeping the other pair off reverses the motor’s rotation direction (e.g., counterclockwise).




%\begin{tabular}{|c|c|c|}
%	\hline
%	LOW (0)	& Low (0)  & Motor OFF\\
%	\hline
%	High (1) & Low (0) & Clockwise \\
%	\hline
%	Low (0) & High (1)  & Counterclockwise \\
%	\hline
%	High (1) & High (1) & Motor OFF \\
%	\hline
%\end{tabular}

\section{\textcolor{red}{Power management}}

\section{\textcolor{red}{Software for Motor control}}
\subsection{\textcolor{red}{Libraries}}
\subsection{\textcolor{red}{Calibration}}

\section{\textcolor{red}{Example Application Motor}}
\subsection{\textcolor{red}{Hardware used in Motor Application}}
\subsubsection{\textcolor{red}{L298N}}
\subsection{\textcolor{red}{Application Discribtion}}

